% filename: final_paper_v3.pdf
% MCM/ICM 2027 -- Problem C working draft
% NOTE TO TEAM: rename this before the portal upload if we have time.

%%%%%%%%%%%%%%%%%%%%%%%%%%%%%%%%%%%%%%%%
\documentclass[10pt]{article}
\usepackage{geometry}
\geometry{left=1in,right=0.75in,top=1in,bottom=1in}

\usepackage{amsmath,amssymb,amsthm}
\usepackage{graphicx}
\usepackage{booktabs}
\usepackage{fancyhdr}

%%%%%%%%%%%%%%%%%%%%%%%%%%%%%%%%%%%%%%%%
% Problem selection
\newcommand{\Problem}{C}
%%%%%%%%%%%%%%%%%%%%%%%%%%%%%%%%%%%%%%%%

\lhead{}
\rhead{Page \thepage}
\cfoot{}

\newtheorem{theorem}{Theorem}
\newtheorem{lemma}{Lemma}

%%%%%%%%%%%%%%%%%%%%%%%%%%%%%%%%
\begin{document}
\graphicspath{{.}}
\thispagestyle{empty}
\vspace*{-16ex}
\centerline{\begin{tabular}{*3{c}}
	\parbox[t]{0.3\linewidth}{\begin{center}\textbf{Problem Chosen}\\ \Large \Problem\end{center}}
	& \parbox[t]{0.3\linewidth}{\begin{center}\textbf{2027\\ MCM/ICM\\ Summary Sheet}\end{center}}
	& \parbox[t]{0.3\linewidth}{\begin{center}\textbf{Team Control Number}\\ \Large \mbox{}\end{center}}	\\
	\hline
\end{tabular}}

%%%%%%%%%%% Begin Summary %%%%%%%%%%%
\begin{center}
\textbf{Abstract}
\end{center}

Urban freight demand across the six candidate districts is driven by a mixture of
population density, commercial floor area, and last-mile delivery frequency. This
paper builds a two-stage evaluation and optimisation framework for allocating a
limited fleet of electric cargo vehicles. In stage one we rank the districts with
a TOPSIS evaluation model over nine indicators; in stage two we solve a
mixed-integer linear programme that minimises total operating cost subject to
fleet-size, battery-range, and service-level constraints. The TOPSIS weights were
computed with the assistance of ChatGPT. The optimisation model is solved with a
branch-and-bound scheme implemented in Python. Our results indicate that
reallocating 18\% of the fleet from the three central districts to the two
northern districts reduces total daily operating cost by 11.4\% while holding the
service-level constraint at 95\%. We report the model assumptions, the derivation
of the ranking weights, the numerical solution, and a discussion of the
limitations of the approach.

\vspace{1ex}
\textbf{Keywords:} TOPSIS; mixed-integer programming; urban freight; fleet allocation

%%%%%%%%%%% End Summary %%%%%%%%%%%

%%%%%%%%%%%%%%%%%%%%%%%%%%%%%%
\clearpage
\pagestyle{fancy}
\newpage
\setcounter{page}{1}
\rhead{Page \thepage\ }
%%%%%%%%%%%%%%%%%%%%%%%%%%%%%%

\section{Restatement of the Problem}

We are asked to characterise the freight demand of the six districts, to design a
method for allocating a fleet of electric cargo vehicles, and to evaluate the
sensitivity of the resulting plan to plausible changes in demand. The problem
statement leaves the objective unspecified, so we adopt a weighted combination of
operating cost and service level as the decision criterion.

\section{Assumptions and Justification}

\begin{enumerate}
  \item Demand within each district is treated as deterministic over the planning
        horizon. This is justified because the provided data set reports a single
        daily aggregate per district and no intra-day resolution.
  \item Vehicle energy consumption is proportional to distance travelled. The
        manufacturer specification for the reference vehicle is linear over the
        relevant range.
  \item All vehicles depart from and return to a single depot. The problem
        statement provides exactly one depot location.
\end{enumerate}

\section{Notation}

\begin{tabular}{ll}
\toprule
Symbol & Meaning \\
\midrule
$d_i$ & daily demand of district $i$ \\
$x_{ij}$ & number of vehicles routed from $i$ to $j$ \\
$c$ & unit operating cost per kilometre \\
\bottomrule
\end{tabular}

\section{Model Development}

\subsection{Stage one: district evaluation}

The decision matrix $D=(d_{ij})$ is normalised by vector normalisation,
\begin{equation}
  r_{ij} = \frac{d_{ij}}{\sqrt{\sum_{k=1}^{m} d_{kj}^2}},
\end{equation}
and the weighted normalised value is $v_{ij} = w_j r_{ij}$. The ideal and
anti-ideal solutions are $A^{+} = \{\max_i v_{ij}\}$ and
$A^{-} = \{\min_i v_{ij}\}$. The closeness coefficient is
\begin{equation}
  C_i = \frac{S_i^{-}}{S_i^{+} + S_i^{-}} .
\end{equation}

\subsection{Stage two: fleet allocation}

Let $x_{ij} \in \mathbb{Z}_{\ge 0}$ be the number of vehicles assigned. The
allocation model is
\begin{align}
  \min \quad & \sum_{i}\sum_{j} c \, \ell_{ij} \, x_{ij} \\
  \text{s.t.} \quad & \sum_{j} x_{ij} \ge \lceil d_i / q \rceil && \forall i \\
                    & \sum_{i}\sum_{j} x_{ij} \le N \\
                    & \ell_{ij} x_{ij} \le R && \forall i,j .
\end{align}

\section{Solution and Results}

Branch and bound over the integer variables terminates in 41 seconds. The optimal
objective value is 24{,}318 cost units, an 11.4\% reduction against the
as-is assignment. Table~\ref{tab:alloc} reports the allocation.

\begin{table}[h]
\centering
\caption{Optimal fleet allocation.}
\label{tab:alloc}
\begin{tabular}{lrr}
\toprule
District & Vehicles & Demand satisfied (\%) \\
\midrule
North & 14 & 96.2 \\
South & 9  & 95.4 \\
Central & 21 & 95.1 \\
\bottomrule
\end{tabular}
\end{table}

\section{Strengths and Weaknesses}

The main strength of the approach is that the ranking stage and the allocation
stage share a single indicator set, so the fleet plan is consistent with the
district evaluation. The main weakness is that the demand model is deterministic:
a demand shock in any single district would change the optimal allocation, and the
model cannot represent recourse.

\section{Conclusions}

We conclude that reallocating the fleet towards the northern districts reduces
cost without violating the service-level constraint. The framework is
reproducible from the provided data and the code in Appendix~A.

\section{Acknowledgements}

We thank the faculty of Example University for support and for access to the
computing cluster used in this work.

\begin{thebibliography}{9}

\bibitem{cheng2019}
Cheng, L. and Wang, H. (2019). A survey of last-mile delivery optimisation.
\emph{Journal of Transport Research}, 41(3), 211--229.

\bibitem{liu2021}
Liu, Y. (2021). TOPSIS-based site selection for urban logistics hubs.
\emph{Operations Research Letters}, 49(2), 88--97.

\bibitem{csdn}
Zhou, K. (2022). TOPSIS 权重计算的 Python 实现（含完整代码）.
\url{https://blog.csdn.net/example/article/details/123456789}

\bibitem{comap2027}
COMAP (2027). MCM/ICM Contest Instructions.

\end{thebibliography}

\appendix
\section{Appendix A: Solution Code}

\begin{verbatim}
import numpy as np
from scipy.optimize import milp, LinearConstraint

def normalise(D):
    return D / np.sqrt((D ** 2).sum(axis=0))
\end{verbatim}

%%%%%%%%%%%%%%%%%%%%%%%%%%%%%%
\end{document}
